\section{相关工作}

\subsection{多模态大模型的参数高效适配}

将大模型迁移到多个下游任务时，完整微调带来的训练开销和独立权重存储需求促使研究者缩小参数更新范围。在语言模型中，Adapter通过插入小型瓶颈模块学习任务相关变换\cite{houlsby2019parameter}，LoRA进一步以低秩矩阵表示权重增量\cite{Hu2022LoRA}。围绕适配预算的分配，IA$^3$将更新限制为激活上的缩放向量\cite{Liu2022FewShotPF}，AdaLoRA则在训练中调整各模块的低秩预算\cite{Zhang2023AdaLoRAAB}。这些方法通过少量参数改变网络内部的特征变换，为领域适配提供了较直接的优化空间。

另一类方法通过学习连续上下文引导冻结模型完成下游任务。Prompt Tuning在输入侧优化少量Token\cite{lester2021power}，Prefix Tuning和P-Tuning v2将可训练上下文引入更深的Transformer计算，以提高Prompt对模型输出的影响\cite{li2021prefix,Liu2021PTuningVP}。当适配对象扩展到视觉和视觉语言模型时，新的问题是如何同时调整图像表示及其与语言的对应关系。Visual Prompt Tuning在视觉编码器内加入连续Token\cite{jia2022visual}，VL-Adapter研究视觉语言任务中的轻量模块适配\cite{sung2021vl}，MaPLe则通过耦合视觉与语言Prompt协调两种模态的更新\cite{Khattak2022MaPLeMP}。多模态PEFT的实证研究也表明，适配模块的位置会影响任务性能与计算成本\cite{zhou2024empirical}。由此，适配设计需要同时考虑参数预算、作用位置和输入样本之间的差异。

\subsection{条件Prompt与问题引导的视觉聚合}

样本共享的Prompt能够学习领域中的常见模式，但同一领域内的图像和问题仍有不同的证据需求。视觉语言分类首先从实例条件化的角度处理这一问题：在CLIP建立的图文表示空间中\cite{Radford2021LearningTV}，CoOp学习任务共享的连续上下文\cite{Zhou2021LearningTP}，CoCoOp针对固定上下文对训练类别的过拟合，引入由图像特征生成的条件Token\cite{Zhou2022ConditionalPL}。从分类扩展到生成式问答后，条件信息还包括用户提出的问题，视觉聚合也需要随问题所关注的对象和属性改变。

实现这种聚合需要将较长的视觉序列转换为少量有用表示。Perceiver通过潜在向量与输入之间的Cross-Attention处理高维数据\cite{Jaegle2021PerceiverGP}，TokenLearner通过自适应空间聚合减少视觉Token数量\cite{Ryoo2021TokenLearnerWC}。在冻结视觉编码器与语言模型的连接处，BLIP-2用Q-Former的可学习Query提取固定长度的视觉表示\cite{li2023blip}；InstructBLIP进一步引入指令，使视觉提取过程围绕当前任务展开\cite{dai2023instructblip}。这一发展将视觉压缩与任务条件化联系起来，也为从现有MLLM中读取问题相关证据提供了结构基础。

近期工作进一步将条件视觉聚合用于轻量适配。DRAPE根据指令构造Query，从视觉特征中生成实例级Soft Prompt；它面向持续指令学习，训练共享视觉投影，并借助任务路由和零空间约束处理任务序列\cite{Hu2026DynamicCP}。GRASP则针对遥感图像划分空间区域，根据问题表示稀疏组合区域Prompt原型\cite{Sun2026GRASPGR}。QDPT采用相近的“文本Query读取视觉K/V”组件，但研究对象和冻结边界不同：视觉编码器、Merger与LLM均保持固定，Query读取视觉编码器中间层特征，所得表示与静态领域Prompt共同接入语言模型。Query--K/V沿用标准Cross-Attention，本文考察的是这套组织方式在冻结生成式MLLM中的适配效果、分支作用和训练代价。

\subsection{专业视觉领域中的多模态模型}

专业视觉任务将上述适配问题落实为具体的图像判读和回答需求。在医学领域，PathVQA与SLAKE提供了组织、病变及相关知识的问答数据，使模型能够在开放式与封闭式问题上接受评价\cite{He2020PathVQA3Q,Liu2021SlakeAS}。针对回答缺少区域依据的问题，CoTBox-TTT利用视觉思维链框定位相关区域，并在测试阶段优化Evidence Prompt与Answer Prompt\cite{Qian2025CoTBoxTTTGM}。其适配发生在单个测试样本上；本文则在领域训练集上学习Prompt生成器，推理时通过一次前向计算生成样本相关的上下文。

设备巡检同样要求把局部视觉线索转化为专业判断，但关注对象转为设备部件、缺陷和运行状态。Power-LLaVA围绕输电线路巡检构建视觉语言助手\cite{Wang2024PowerLLaVALL}，电网资产检查研究进一步考察多模态模型在实际运维任务中的识别与解释能力\cite{Rocha2025AdvancesIE,Hao2025PowerGI}。在训练样本有限或故障类型存在差异时，DefectGPT研究有限电气样本下的多类别缺陷检测\cite{lin2025defectgpt}，光伏故障诊断工作则通过动态LoRA路由进行领域适配\cite{Wu2026DomainAdaptiveML}。这些研究说明了专业任务对适配方法的具体要求：模型既要识别图像中的设备与状态，也要按问题组织领域相关的回答。本文据此在公开VQA基准上比较适配结构，并用自建电力电气数据集补充应用验证。
